% ============================================================================
% Plantilla de tesis de MAESTRÍA — Normas APA, sexta edición (español)
% Fuente de reglas: Centro de Escritura Javeriano (2018), Normas APA 6ª ed.
% Motor: pdfLaTeX (fuente Times vía newtx). Bibliografía: biblatex + biber.
%
% Compilar:
%   pdflatex <archivo> && biber <archivo> && pdflatex <archivo> && pdflatex <archivo>
% ============================================================================

\documentclass[letterpaper,12pt]{article}   % §1.1 papel carta · §1.3 12 pt

% ---------------------------------------------------------------------------
% Idioma y codificación
% ---------------------------------------------------------------------------
\usepackage[spanish,es-nodecimaldot]{babel}
\usepackage[utf8]{inputenc}
\usepackage[T1]{fontenc}

% ---------------------------------------------------------------------------
% Fuente: Times New Roman (clon newtx) — §1.3
% ---------------------------------------------------------------------------
\usepackage{newtxtext,newtxmath}

% ---------------------------------------------------------------------------
% Márgenes 2.54 cm en los cuatro lados — §1.2
% ---------------------------------------------------------------------------
\usepackage[letterpaper,margin=2.54cm]{geometry}

% ---------------------------------------------------------------------------
% Interlineado 2.0 — §1.4
% ---------------------------------------------------------------------------
\usepackage{setspace}
\doublespacing

% ---------------------------------------------------------------------------
% Alineación a la izquierda, SIN justificar — §1.4
% ---------------------------------------------------------------------------
\usepackage[document]{ragged2e}
\setlength{\RaggedRightParindent}{1.27cm}   % conserva sangría con ragged

% Sangría 1ª línea (5 espacios ≈ 0.5 in) — §1.5 · sin espacio entre párrafos — §1.4
\setlength{\parindent}{1.27cm}
\setlength{\parskip}{0pt}

% ---------------------------------------------------------------------------
% Gráficos, tablas (solo líneas horizontales — §1.6, §3.3) y comillas
% ---------------------------------------------------------------------------
\usepackage{graphicx}
\graphicspath{ {./images/} }
\usepackage{booktabs}     % \toprule \midrule \bottomrule — NO usar \hline por celda
\usepackage{csquotes}

% ---------------------------------------------------------------------------
% Encabezados APA de 5 niveles — §2.2
%   N1: centrado, negrita
%   N2: alineado izq., negrita
%   N3: sangría, negrita, run-in, termina en punto
%   N4: sangría, negrita+cursiva, run-in, termina en punto
%   N5: sangría, cursiva, run-in, termina en punto
% Recuerda: mayúscula inicial solo en la 1ª palabra — §2.1
% ---------------------------------------------------------------------------
\usepackage{titlesec}

% article: section=N1, subsection=N2, subsubsection=N3, paragraph=N4, subparagraph=N5
\titleformat{\section}{\normalfont\bfseries\centering}{\thesection}{1em}{}
\titleformat{\subsection}{\normalfont\bfseries\raggedright}{\thesubsection}{1em}{}
\titleformat{\subsubsection}[runin]{\normalfont\bfseries}{}{0pt}{}[.\quad]
\titleformat{\paragraph}[runin]{\normalfont\bfseries\itshape}{}{0pt}{}[.\quad]
\titleformat{\subparagraph}[runin]{\normalfont\itshape}{}{0pt}{}[.\quad]

% Sangría de los niveles run-in (N3, N4, N5)
\titlespacing{\subsubsection}{1.27cm}{\parskip}{0.5em}
\titlespacing{\paragraph}{1.27cm}{\parskip}{0.5em}
\titlespacing{\subparagraph}{1.27cm}{\parskip}{0.5em}

% ---------------------------------------------------------------------------
% Numeración de páginas (esquina superior derecha)
% ---------------------------------------------------------------------------
\usepackage{fancyhdr}
\setlength{\headheight}{14.5pt}
\pagestyle{fancy}
\fancyhf{}
\fancyhead[R]{\thepage}
\renewcommand{\headrulewidth}{0pt}

% ---------------------------------------------------------------------------
% Bibliografía APA — §12 (sangría francesa y orden alfabético automáticos)
% ---------------------------------------------------------------------------
\usepackage[backend=biber,style=apa,language=spanish,sortcites=true,sorting=nyt]{biblatex}
\DeclareLanguageMapping{spanish}{spanish-apa}
\addbibresource{referencias.bib}

% ---------------------------------------------------------------------------
% Hiperenlaces (en negro para impresión formal)
% ---------------------------------------------------------------------------
\usepackage[colorlinks=true,linkcolor=black,citecolor=black,urlcolor=black,plainpages=false,hypertexnames=false]{hyperref}

% ============================================================================
\begin{document}

% ---------------------------------------------------------------------------
% PORTADA de maestría
% NOTA: APA 6ª ed. no define portada (guía §23). Ajusta los datos de tu
% universidad/programa; conserva "contenido-titulo" (el script lo sustituye).
% ---------------------------------------------------------------------------
\begin{titlepage}
    \centering
    \thispagestyle{empty}
    \singlespacing
    \vspace*{1cm}

    {\bfseries UNIVERSIDAD [NOMBRE DE LA UNIVERSIDAD]\par}
    {\bfseries FACULTAD DE [NOMBRE DE LA FACULTAD]\par}
    {[Maestría en \_\_\_\_\_\_\_\_\_\_\_\_]\par}

    \vfill

    {\LARGE\bfseries contenido-titulo\par}

    \vspace{1cm}
    {[Trabajo de grado para optar al título de Magíster en \_\_\_\_\_\_\_\_\_\_]\par}

    \vfill

    {\bfseries Autor\par}
    {[Nombre completo del autor]\par}

    \vspace{0.8cm}
    {\bfseries Asesor\par}
    {[Nombre completo del asesor]\par}

    \vfill

    {[Ciudad, País]\par}
    {[Año]\par}
\end{titlepage}

% ---------------------------------------------------------------------------
% Preliminares (numeración romana). Descomenta lo que necesites.
% ---------------------------------------------------------------------------
\pagenumbering{roman}
%   \include{src/dedicatoria}
%   \include{src/agradecimientos}
%   \include{src/resumen}
    \tableofcontents
    \newpage
%   \listoffigures
%   \listoftables

% ---------------------------------------------------------------------------
% Cuerpo del documento (numeración arábiga)
% ---------------------------------------------------------------------------
\clearpage
\pagenumbering{arabic}

% Escribe cada apartado en src/ y descoméntalo aquí.
% Cada \section es un encabezado de Nivel 1 (mayúscula solo inicial — §2.1).
    \include{src/introduccion}
%   \include{src/marco_teorico}
%   \include{src/metodologia}
%   \include{src/resultados}
%   \include{src/discusion}
%   \include{src/conclusiones}

% ---------------------------------------------------------------------------
% Referencias — §12 (biblatex-apa: sangría francesa + orden alfabético)
% ---------------------------------------------------------------------------
\newpage
\printbibliography[title=Referencias]

% ---------------------------------------------------------------------------
% Anexos (opcional)
% ---------------------------------------------------------------------------
%   \newpage
%   \appendix
%   \include{src/anexos}

\end{document}

% ============================================================================
% RECORDATORIOS APA que el motor NO automatiza (revisa al escribir):
%
%  Citas en el texto (biblatex-apa) — §5–§11:
%    \textcite{clave}   -> Apellido (Año) afirma...           (basada en autor)
%    \parencite{clave}  -> ...idea... (Apellido, Año)         (basada en texto)
%    \parencite[p.~90]{clave} -> ... (Apellido, Año, p. 90)   (cita textual — §6)
%    3–5 autores: 1ª cita todos, luego "et al." (biblatex lo hace) — §10.2
%    Autor corporativo, s.f., "como se citó en", etc. — §11
%
%  Cita textual < 40 palabras: entre comillas, en el párrafo, sin cursiva — §7
%  Cita textual > 40 palabras: bloque con sangría, sin comillas — §8
%    \begin{quote} ... texto ... \parencite[p.~##]{clave} \end{quote}
%
%  Títulos: mayúscula SOLO en la primera palabra, no en mayúscula sostenida — §2.1
%
%  Tablas — §3: número + título breve en cursiva ARRIBA; sin líneas internas
%    (solo \toprule/\midrule/\bottomrule); nota debajo; si es propia: "Autoría propia."
%
%  Figuras — §4: "Figura X." en cursiva; nota debajo con el mismo tamaño de letra.
% ============================================================================
